The preceding sections concern how a document is represented. This one concerns
what is selected once candidates are retrieved, which is where the study's second
claim lives: that the right amount of diversity depends on who consumes the
results. This section settles the measurable half of that claim, namely what each
policy actually delivers; the consumer half is future work
(Section~\ref{sec:limitations}).

\paragraph{Setting.}
Each of 120 constructed queries is associated with several distinct facts, and
each candidate document realises at most one of them, so the subtopics a
selection covers can be counted exactly rather than inferred. Deriving subtopics
by clustering would have been circular, since one of the policies under test
selects by cluster. Candidates are gated to the upper half by relevance first, as
a first pass would do, and each policy then selects ten items.

\begin{table}[t]
\caption{Selection policies on a constructed coverage task. All coverage gains
over relevance ranking are significant under Holm-corrected paired permutation
tests. Relevance differences are within noise, an artefact of the strong gate
discussed below.}
\label{tab:diversity}
\centering
\begin{tabular}{lrrrr}
\toprule
policy & S-recall & $\alpha$-nDCG & nDCG@10 & answerless rate \\
\midrule
relevance only          & 0.628 & 0.692 & 0.965 & 0.033 \\
MMR, $\lambda=0.7$      & 0.666 & 0.723 & 0.964 & 0.032 \\
MMR, $\lambda=0.3$      & 0.818 & 0.840 & 0.963 & 0.036 \\
DPP                     & \textbf{0.879} & \textbf{0.877} & 0.963 & 0.035 \\
cluster round robin     & 0.666 & 0.724 & 0.967 & 0.028 \\
outlier harvest, 20\%   & 0.716 & 0.734 & 0.966 & 0.031 \\
outlier harvest, 40\%   & 0.772 & 0.775 & 0.967 & 0.028 \\
\bottomrule
\end{tabular}
\end{table}

\paragraph{The operator we proposed is not the best operator.}
Table~\ref{tab:diversity} is not the result we expected. Harvesting the
low-similarity tail of a gated pool does raise coverage substantially, by 0.144
S-recall at a 40 percent quota, but a determinantal point process raises it by
0.251 and maximal marginal relevance at $\lambda=0.3$ by 0.190. We report this
because it sharpens rather than weakens the position we are arguing for. The
claim worth defending is not that a particular diversification operator is best;
it is that the quantity of diversity is a function of the consumer, and that
function applies to whichever operator a system already has. On this evidence the
strongest implementation of the idea is a determinantal point process whose
quality-diversity balance is set per consumer, not the harvest rule.

\paragraph{Under a strong gate, coverage is nearly free.}
No policy paid a measurable relevance penalty, and the answerless rate did not
move. Both follow from gating: when the pool is already bounded for relevance,
spending part of the budget away from its centre costs little and risks little.
This is an argument for the gate rather than evidence that diversity is free, and
we would expect the trade to reappear with a weaker filter or a stricter
relevance criterion.
